\subsection{Итоговый проект: поиск и спасение}
\label{sec:swarm-project}

Цель раздела: собрать группу из двух дронов с разными ролями: разведчик ищет, спасатель ждёт сигнал и прилетает в найденную точку.

Протокол одно сообщение: \verb|'НАЙДЕН x y'| --- координаты найденного пострадавшего. Спасатель стоит на земле и ждёт его, разведчик летит «змейкой» над площадкой $3\times 3$~м; как только проходит рядом с пострадавшим, рассылает рою координаты и садится. Запускают сначала спасателя, затем разведчика.

\begin{codefile}[python]{python-onboard/swarm/swarm_09_rescue_rescuer.py}
\end{codefile}

\begin{codefile}[python]{python-onboard/swarm/swarm_09_rescue_searcher.py}
\end{codefile}

\paragraph{Что происходит (спасатель).} Минуту ждёт письмо (опрос ящика с паузой, как раньше). Разбор строго проверяет формат: \verb|text.split()| должно дать три части, первая --- \verb|'НАЙДЕН'|, иначе сообщение игнорируется. Две оставшиеся части преобразуются в \verb|float|: это цель. Дальше --- обычный полёт: \verb|arm|, \verb|takeoff|, точка цели с ожиданием \verb|point_reached()|, зелёный светодиод, пауза и посадка. Если сигнала не было, спасатель так и стоит на земле и закрывает соединение.

\paragraph{Что происходит (разведчик).} Маршрут «змейка» --- список из восьми точек. Пока дрон летит к очередной точке, цикл ожидания прибытия каждые 0.1~с измеряет расстояние до пострадавшего (\verb|math.hypot|) по текущей позиции; если оно меньше \verb|SEEN_RADIUS| (0.8~м), разведчик запоминает свою позицию как место находки и выходит из обоих циклов. В примере нет камеры: «видит» дрон по известной точке \verb|VICTIM|, чтобы проверка была повторяемой. Найдя цель, он рассылает всем соседям \verb|'НАЙДЕН x y'| и ждёт восемь секунд, пока спасатель долетит, затем садится.

\begin{warnbox}[Безопасность группы]
Два дрона в одном небе --- уже риск столкновения: в примере нет защиты от встречного движения, а спасатель летит в точку, которую назвал разведчик. Разводите дроны по высотам (разведчик 1.5~м, спасатель 1~м), не посылайте их через одну точку одновременно и не запускайте группу рядом с людьми. Протокол без подтверждений может потерять сообщение: в реальном проекте спасатель отвечает «принял», а разведчик повторяет сигнал, пока не получит ответ.
\end{warnbox}

\begin{taskbox}[Найти и вызвать]
Мини-проект. Доработайте протокол: спасатель отвечает «принял x y»; разведчик повторяет «НАЙДЕН» раз в 2~секунды до подтверждения (не больше пяти раз) и печатает результат. Добавьте третьего дрона, который слушает рассылку и зажигает синий свет как «наблюдатель».
\expected{при подтверждении разведчик печатает «Спасатель принял»; без спасателя --- «Подтверждения нет» после пяти попыток и посадка; наблюдатель включает синий свет при первом сообщении.}
\mission{Мини~2 SDK2}{64}
\end{taskbox}

\paragraph{Итоговая миссия трека.} Миссия 64 --- звено командира и ведомого: «Звено готово: N из N» печатается до \verb|arm|, затем командир летит маршрут $(1,0,1) \to (1,1,1) \to (0,1,1)$, а ведомый держится позади; диод 0 командира --- жёлтый, зелёный, выключен. Основа --- примеры о синхронном взлёте и следовании. Вы можете собрать такую связку из программ этой главы.

\subsection{Проверьте себя}

\begin{enumerate}
  \item Чем рой Swarm Kit отличается от схемы «все дроны подчиняются наземной станции»?
    \begin{enumerate}[label=(\alph*)] \item ничем \item юниты равноправны, центрального сервера нет \item в рое есть главный юнит \end{enumerate}
    Ответ: (b). Все юниты подключены к одной сети и общаются напрямую.
  \item Что такое Heartbeat?
    \begin{enumerate}[label=(\alph*)] \item периодический сигнал «я жив» с типом юнита и каталогом данных \item письмо в ящик \item поток телеметрии \end{enumerate}
    Ответ: (a). Из Heartbeat узнают, кто в сети и на что можно подписаться.
  \item Что вернёт \verb|get_unit_list()|, если рядом работает один дрон кроме вашего?
    \begin{enumerate}[label=(\alph*)] \item два идентификатора \item список из одного идентификатора соседа \item пустой список \end{enumerate}
    Ответ: (b). Свой \texttt{machine\_id} в список не входит.
  \item \verb|send_message_to_unit('pioneer9', 'привет')| вернул \verb|False|. Почему?
    \begin{enumerate}[label=(\alph*)] \item юнита \texttt{pioneer9} нет в списке соседей \item ящик переполнен \item сообщение слишком длинное \end{enumerate}
    Ответ: (a). \verb|True| означает «адресат найден, сообщение ушло», но не «прочитано».
  \item Что вернёт \verb|get_next_message()|, если ящик пуст?
    \begin{enumerate}[label=(\alph*)] \item подождёт и вернёт сообщение \item \verb|None| \item пустую строку \end{enumerate}
    Ответ: (b). Метод не ждёт; ждать нужно опросом с паузой и пределом.
  \item Почему цикл ожидания сообщения должен иметь паузу и предел?
    \begin{enumerate}[label=(\alph*)] \item пауза не нагружает процессор, предел спасает от вечного ожидания \item иначе сообщения теряются \item так требует \texttt{swarm-server} \end{enumerate}
    Ответ: (a).
  \item Что случится при \verb|get_telemetry_from_unit(id, timeout_sec=0)|?
    \begin{enumerate}[label=(\alph*)] \item подписка и возврат накопленного, возможно пустого словаря \item ожидание без предела \item исключение \end{enumerate}
    Ответ: (a). \verb|None| вместо числа --- ожидание всех потоков без ограничения.
  \item Ведомый использует локальные координаты ведущего. При каком условии это верно?
    \begin{enumerate}[label=(\alph*)] \item всегда \item у дронов общая система координат \item когда дроны одинаковые \end{enumerate}
    Ответ: (b). Иначе ведомый получит чужие координаты в своей системе.
  \item Почему команду «взлёт» посылают только после ответов «готов» от всех?
    \begin{enumerate}[label=(\alph*)] \item иначе часть группы не взлетит или взлетит не вовремя \item так быстрее \item это требует библиотека \end{enumerate}
    Ответ: (a). Подтверждения превращают рассылку в согласованное действие; отказ любого участника отменяет взлёт.
  \item Какие ограничения у письма в симуляторе, о которых нельзя забывать в настоящей сети?
    \begin{enumerate}[label=(\alph*)] \item доставка мгновенная и без потерь только в симуляторе \item письма шифруются \item ящик бесконечен \end{enumerate}
    Ответ: (a). В симуляторе нет задержек и потерь; ящик ограничен 100 сообщениями; подтверждений доставки нет.
\end{enumerate}

\subsection{Задания к главе}

\begin{taskbox}[Перекличка перед полётом]
Командир перед взлётом проверяет состав: если число соседей меньше ожидаемого, печатает, кого не хватает, и не взлетает. Затем рассылает «готовность» и ждёт ответы с пределом 15~секунд.
\expected{при полном составе все отвечают и командир взлетает; при нехватке дрона печатается его идентификатор, и моторы не запускаются.}
\mission{Мини~2 SDK2}{57, 63}
\end{taskbox}

\begin{taskbox}[Цепочка ведомых]
Мини-проект. Три дрона летят цепочкой: второй следует за первым, третий --- за вторым на метр позади. Для этого каждый читает телеметрию своего предшественника, а смещения и номера не зашиты в код, а приходят командой от первого.
\expected{в 3D три дрона держатся друг за другом без столкновений; при пропаже предшественника ведомый зависает и садится.}
\end{taskbox}

\subsection{Миссии симулятора по разделам главы}

Главы 18--19 трека \ui{Мини 2 · SDK2}:

{\small
\begin{longtable}{@{}>{\raggedright\arraybackslash}p{0.09\linewidth}>{\raggedright\arraybackslash}p{0.38\linewidth}>{\raggedright\arraybackslash}p{0.45\linewidth}@{}}
\toprule
\textbf{№} & \textbf{Тема миссий} & \textbf{Раздел главы} \\
\midrule
\endhead
57 & Кто в рое (свой \texttt{machine\_id}, соседи) & кто в сети \\
58 & Сообщение соседу, доставка и отказ & сообщения \\
59 & Ping-pong с соседом & сообщения: ping-pong \\
60 & Рассылка всем соседям & рассылка \\
61 & Телеметрия соседа (высота, состояние) & телеметрия соседей \\
62 & Лететь за ведущим & ведущий и ведомый \\
63 & Взлёт по готовности & синхронный взлёт \\
64 & Итог: звено командира и ведомого & итоговый проект \\
\bottomrule
\end{longtable}
}

\subsection*{Итоги}
\begin{itemize}
    \item Рой Swarm Kit одноранговый: юниты равноправны, у каждого есть \texttt{machine\_id}; центрального сервера нет.
    \item Три способа общения: Heartbeat («кто в сети» и каталог данных), телеметрия по подписке, сообщения в почтовый ящик юнита.
    \item \verb|PioneerSwarmUnit()| подключает программу; \verb|get_unit_list()| --- соседи без себя; \verb|close()| --- в конце; свойства \verb|.id| и \verb|.unit_type|.
    \item Сообщение --- строка; \verb|send_message_to_unit| возвращает \verb|True| или \verb|False| (юнита нет); \verb|get_next_message()| не ждёт --- опрос с паузой и пределом.
    \item \verb|get_telemetry_from_unit(id, timeout_sec)| возвращает \verb|(тип, словарь)| или \verb|None|; словарь может быть неполным; значения --- сообщения ROS~2.
    \item Следование и строй берут позицию соседа из телеметрии; это верно только при общей системе координат, а защиты от столкновений нет.
    \item Групповые действия строят на подтверждениях: «готовность» $\to$ «готов» $\to$ «взлёт»; протокол из строк разбирают с проверкой формата.
    \item Все программы проверены на симуляторе с мгновенной сетью; на настоящих дронах авторы их не запускали, надёжность доставки и порядок сообщений в настоящей сети не гарантированы.
\end{itemize}
